\section*{摘要}
\addcontentsline{toc}{section}{摘要}
本实验以一台 TonyPi 人形机器人和多台 TurtleBot2 为平台，研究抱物穿行、视觉航路控制及多方式跟随的系统集成。任务要求人形机器人抱持物体完成至少一次左转和一次右转，终点放置物体；各移动机器人以可区分的方式保持跟随，并满足避碰、目标保持和终点停靠要求。

人形机器人部分采用“动作组执行＋视觉反馈”的分层结构：以抓取末帧固定上肢姿态，复用下肢步态；先通过开环路线定位动作偏移问题，再实现灰色挡板道路居中，最新版本进一步采用 A/B/C 固定标记进行测距、对中与转角切换。方式二采用赛前 GMapping 建图、比赛 AMCL 定位、AR 观测变换至地图、历史轨迹回退生成目标以及 move\_base/DWA 导航，并设置独立速度门控。报告给出分工模板、接口约定、实现原理、已知验证范围和实验记录表，供各组并行补充。

现有仓库记录支持局部实机运行与部分软件验证，尚不足以证明整条路线及全队比赛验收完成。本文不将仿真结果、默认参数或设计目标作为实测成绩。

\noindent\textbf{关键词：}TonyPi；TurtleBot2；动作组；视觉标记；地图导航；多机协同
